\subsection{A sharp characteristic-prime criterion and the endpoint obstruction}
\label{sec:rh-sp4-curve-frontier}

The finiteness and endpoint premises of
Theorem~\ref{thm:rh-sp4-conditional-length} can be stated
as precise arithmetic exclusions. This is useful because it
separates the genuinely characteristic-prime difficulty
from the characteristic-zero count. No flatness of the
global equation quotient is assumed in the following
reductions.

\begin{theorem}[Equivalent curve tests for the full invariant algebra]
\label{thm:rh-sp4-curve-criterion}
Keep the complete normalized ring \(S\) of the actual RH
rank-four model and its augmentation, with coefficient
uniformizer \(\varpi\) and finite residue field \(k\).
By a characteristic-prime curve we mean any local
\(O\)-algebra map
\[
                  f:S\longrightarrow D
\]
to a complete discrete valuation ring of characteristic
\(\ell\) with finite residue field, using the given
map \(O\to k\to D\). Let \(u_f\) be its original
normalized representation, and let \(\bar r_0\) be the
residual diagonal representation, viewed over the constant
field of \(D\). The following are equivalent:
\begin{enumerate}
\item The full invariant algebra \(S_{(0,0)}\) is finite
as an \(O\)-module (premise (F)).
\item On every characteristic-prime curve,
\(\operatorname{tr}u_f(g)=\operatorname{tr}\bar r_0(g)\)
for every \(g\in G_{\mathbb Q}\).
\item On every such curve the semisimplified representation
over the algebraic closure of \(\operatorname{Frac}D\)
has finite image.
\item There is neither an absolutely irreducible
characteristic-prime curve nor a block-diagonal
characteristic-prime curve whose first block is
absolutely irreducible.
\end{enumerate}
In condition 4 the blocks and all local conditions are
exactly those of the model. In condition 3 no fixed
finite-index subgroup needs to work for all curves.
\end{theorem}

\begin{proof}
First record the local flag that survives specialization
to characteristic \(\ell\).
Over the fraction field of the actual local domain
\(L^\circ\), the original spectral equations give the
ordered invariant full flag and the prescribed inertia
characters. The values of the graded characters at
\(\mathrm{ph}\) are fixed as in (SR1).
Consider the projective symplectic flag scheme over \(L^\circ\)
(the isotropic flag conditions are closed),
and in it impose stability under every local group
matrix, the prescribed scalar actions on its graded
lines at inertia, and the prescribed graded actions
at \(\mathrm{ph}\).
These are closed equations: preservation of a universal
subbundle is expressed by vanishing of its map to the
corresponding quotient bundle, and the scalar action
is the analogous map on each successive quotient.
Arbitrary intersections of these closed subschemes
are closed (the flag scheme is Noetherian).
The resulting scheme is proper over \(L^\circ\)
and its image contains the generic point.
Because \(L^\circ\) is a domain, the image is all of
\(\operatorname{Spec}L^\circ\).
Every field-valued specialization of the local component
thus admits such a flag over an algebraic closure of
the field. This applies to the characteristic-prime
curve after restriction to the original local ring
\(C\), which is a quotient of \(L^\circ\).
It does not assert that the original spectral
denominators are units in the characteristic-prime
fibre, nor that the flag there is unique.

The coordinate-projector proof of
Lemma~\ref{lem:rh-sp4-graph} works over \(D\) as well.
Its only denominators are the unit resultant of the
two block characteristic polynomials at \(\mathrm{ph}\)
and the unit difference of the diagonal entries of
the compressed marked matrices.
In particular, the support graph of \(u_f\) is
transitive, reflexive and self-reflecting.
If a strongly connected component has size two
across the coordinate blocks, form its two-dimensional
coordinate subquotient: take the span of the indices
\(k\) with \(k\rightsquigarrow i\) and delete the
two indices in the component to obtain the invariant
subspace to be quotiented out.
The local flag just proved supplies an invariant
line over an algebraic closure of \(\operatorname{Frac}D\).
On this subquotient \(\mathrm{ph}\) is diagonal with
distinct eigenvalues, so the invariant line is a
coordinate line. One of the two globally nonzero
off-diagonal entries thus vanishes on \(D_\ell\).
Divide that entry by its minimum nonzero valuation
factor, now a power of the uniformizer of \(D\),
and reduce modulo its maximal ideal.
The same two-dimensional multiplication identity
gives the nonzero residual cocycle for an off-block
character used in Lemma~\ref{lem:rh-sp4-graph}.
It is zero locally at \(\ell\), and satisfies the
original away-from-\(\ell\) inertia-kernel condition.
The residual \texttt{XVan} forces it to be a
coboundary and the separation at \(\mathrm{ph}\)
forces its primitive to be zero, a contradiction.
This uses the original residual character, not a
characteristic-zero local cohomology calculation.
Thus the same three support alternatives hold on
curves as in that lemma, but we do not use its
characteristic-zero ordering of the diagonal type.

We next prove a useful trace argument in every
coefficient characteristic occurring here, including
\(\ell=3\).
Put \(F_D=\operatorname{Frac}D\).
The four residual characters in \(\bar r_0\) are
pairwise distinct: the two blocks separate at
\(\mathrm{ph}\), and within a block their ratio
\(\bar\Psi\) or its inverse is nontrivial at \(g_0\).
Thus their diagonal matrix algebra over \(F_D\)
is the full diagonal algebra, by independence of
distinct characters. Its ordinary matrix trace
pairing is nondegenerate, since each coordinate
has rank one.

Assume that the trace equality of condition 2
holds on the given curve.
Let \(\mathscr M\) be the \(F_D\)-algebra spanned
by the group matrices of \(u_f\).
There is a surjective homomorphism
\[
        \mathscr M\longrightarrow F_D^4,\qquad
                 u_f(g)\longmapsto\bar r_0(g).
                                                               \tag{CF1}
\]
To see well-definedness, any group-algebra expression
mapping to zero in \(\mathscr M\) has residual trace
pairing zero against every group element: use the
trace equality at all products of group elements.
The nondegenerate diagonal trace pairing then makes
its image in \(F_D^4\) zero. This argument does not
divide by the dimension of a simple summand.
The map (CF1) preserves the traces.
Every coordinate projector \(E_{ii}\) lies in the
integral representation algebra; its image under
(CF1) is a diagonal idempotent whose reduction is
\(E_{ii}\), by using its integral group-algebra
expression. Its entries are zero or one, so it
is exactly \(E_{ii}\).
If both opposite off-diagonal entries occur
nontrivially, \(\mathscr M\) contains \(E_{ij}\)
and \(E_{ji}\). Both map to zero under (CF1), by
multiplication with the coordinate projectors.
Their product \(E_{ii}\) does not, a contradiction.
The support is consequently diagonal type, and
\[
            u_f(g)_{ii}=\operatorname{tr}(E_{ii}u_f(g))
                      =\bar r_0(g)_{ii}.                 \tag{CF2}
\]
In particular no irreducible block of dimension
at least two can occur under condition 2.

For diagonal-type support on any curve, its four
diagonal entries are global multiplicative characters.
Their inertia values at \(\ell\) have bounded finite
order: the local flag in characteristic \(\ell\)
has the four prescribed finite residual inertia
characters, and the diagonal characters are
eigenvalues of the same inertia matrices.
At the other finite primes the original inertia
kernel has finite index, and outside the given finite
set the representation is unramified.
A single positive power of each diagonal character
is therefore unramified at every finite prime;
class field theory for \(\mathbb Q\) makes that
power trivial. A finite-order element of \(D^\times\)
lies in the constant field, since it is algebraic
over \(k\). Its residual value determines it.
Thus diagonal-type support implies (CF2) without
assuming condition 2.

For completeness a diagonal-type support graph
also makes every balanced nonconstant monomial
in entry displacements vanish on the curve.
Choose rational numbers \(d_i\) strictly decreasing
along the strict arrows of the partial order, so
\(d_i>d_j\) for \(i\rightsquigarrow j\), \(i\ne j\).
Replace \(d_i\) by
\((d_i-d_{3-i})/2\). The inequalities remain strict
by reflection and the new values have the form
\((a,b,-b,-a)\). These are the values of a linear
functional on the two coordinate weights
\(w_i=(\lambda_i,\mu_i)\), up to a common additive
constant: use coefficients \(a-b,a+b\).
Every nonzero off-diagonal entry then has strictly
positive weight for that functional. Their nonempty
products cannot have full weight zero.
The diagonal entry displacements are zero by
(CF2), proving the balanced-monomial assertion.

Now (F) implies condition 2: the invariant trace
elements on a curve are integral over the finite
residue field and reduce to their base values, hence
are those constants. Conversely, under condition 2,
the preceding argument makes every nonconstant
balanced monomial zero on every curve.
We give the exact integral finiteness step, to
account also for the possible nilpotents and torsion.
Choose the finitely many homogeneous generators
of the original power-series presentation of \(S\),
and finitely many nonconstant balanced monomials
generating the zero-weight monoid by Dickson's lemma.
Their images \(b_1,\ldots,b_s\) generate
\(S_{(0,0)}\) as a complete local quotient of a
finite-variable power-series ring over \(O\).
If some \(b_j\notin\sqrt{\varpi S}\), choose a prime
of \(S/\varpi S\) avoiding it. In the resulting
complete local domain its image is nonzero and
noninvertible. Extend that image to a system of
parameters; a minimal prime of the other parameters
avoiding it gives a one-dimensional domain quotient,
by the height theorem. Its finite normalization is
a complete DVR of characteristic \(\ell\) with
finite residue field. The resulting curve does not
kill \(b_j\), a contradiction.
This is the curve argument of
Proposition~\ref{prop:support-integral-finiteness-criterion},
with the same integral rings and now a two-dimensional
weight monoid. Therefore each \(b_j^{n_j}\in\varpi S\).
Weight projection is \(S_{(0,0)}\)-linear and yields
\((\varpi S)\cap S_{(0,0)}=\varpi S_{(0,0)}\).
It follows that a power of the maximal ideal
of the invariant ring lies in \(\varpi S_{(0,0)}\).
Its quotient modulo \(\varpi\) is finite-dimensional
over \(k\), and its maximal-ideal and \(\varpi\)-adic
topologies are equivalent. Complete successive
approximation lifts a finite generating set over
\(k\) to a finite generating set over \(O\).
This proves (F) on the full invariant algebra,
without deleting any component of \(S\).

Condition 3 implies condition 2 because every
eigenvalue of every matrix in a finite-image
semisimplification is a root of unity, hence
algebraic over \(k\); the characteristic polynomial
coefficients, lying in \(D\), are then the unique
constant-field lifts of their residual values.
Conversely condition 2 gives diagonal-type support
by (CF1), and (CF2) makes its semisimplification a
sum of the four original finite residual characters.

Finally, condition 2 excludes both curves in
condition 4 by (CF1). Under condition 4 use the
three support alternatives proved above.
Full support is forbidden. In the block alternative
at least one interblock direction is absent.
Scale the other interblock direction towards zero
by the first torus using powers of the uniformizer
of \(D\) (or its inverse in the direction that makes
the surviving entries smaller).
The finite set of entry-displacement generators
converges to its block-diagonal specialization
and remains in the maximal ideal of \(D\).
Continuous formal evaluation preserves the exact
equation ideal. This produces a characteristic-prime
curve with an irreducible first block, also
forbidden by condition 4. Only diagonal type is
possible, which already gives condition 2.
This completes the equivalences.
\end{proof}

\begin{proposition}[The precise rank-two endpoint obstruction]
\label{prop:rh-sp4-even-endpoints}
For any block-diagonal characteristic-zero point of
the actual model, write \(U,V\) for its two blocks.
At every good arithmetic Frobenius \(\sigma_p\), for
\(p\nmid M\) as in Theorem~\ref{thm:rh-symplectic-spectrum},
and any coefficient isomorphism \(\iota\), one has
\[
 |\iota\det U(\sigma_p)|=p^{e+2},\qquad
 |\iota\det V(\sigma_p)|=p^e .                            \tag{CF3}
\]
Thus its four eigenvalues cannot all lie in a strict
common real-power interval
\[
             p^{c-1/2}<|\iota\lambda|<p^{c+1/2}
                                                               \tag{CF4}
\]
for any real \(c\).
In particular premise (E) would follow from
such a single good-prime strict interval for
every block-diagonal point with irreducible first
block, even if the prime and \(c\) depend on the point.

The first block, twisted by \(\alpha^{-1}\), is an
ordinary two-dimensional even representation
with distinct Hodge--Tate weights \(-e,0\) and
split residual representation
\(\bar\Psi\oplus1\).
If the first block is absolutely irreducible,
the twisted representation remains so. Its
determinant is a finite-order twist of
\(\varepsilon^e\).
The analogous block on a characteristic-prime
curve, twisted by \(\bar\alpha^{-1}\), is
locally triangular, has finite-order determinant,
and has the same split residual representation.
It cannot have finite image if absolutely irreducible.
\end{proposition}

\begin{proof}
The characteristic-zero local flag has the
labeled values (SR1). Its intersection with the
two blocks, distinguished by the disjoint
spectra at \(\mathrm{ph}\), gives the two
graded inertia characters \(\tau,\alpha\)
in the first block and \(\Psi,1\) in the
second. Thus the determinant characters
have cyclotomic exponents \(e+2,e\) on
inertia at \(\ell\), respectively, times
finite-order factors.
Away from \(\ell\) their inertia images are finite
by the original kernel condition.
Dividing by the indicated cyclotomic powers
therefore gives finite-order global characters,
as in the class-field argument above.
At the specified good prime their absolute
values are one under every coefficient
isomorphism. The cyclotomic value at
arithmetic Frobenius is \(p\), proving (CF3).

If every eigenvalue satisfies (CF4), the two
determinant magnitudes lie strictly between
\(p^{2c-1}\) and \(p^{2c+1}\).
Their ratio is strictly less than \(p^2\),
contradicting (CF3). In particular even a
strict half-power purity window suffices;
one need not prove an exact common Weil weight
or an automorphic identification for these
endpoints.

The residual first block is
\(\operatorname{diag}(\bar\tau,\bar\alpha)\);
division by \(\bar\alpha\) gives
\(\bar\Psi\oplus1\). The local inertia
characters after division by \(\alpha\)
are \(\Psi,1\) and hence have Hodge--Tate
weights \(-e,0\).
The block is potentially crystalline by the
local finite-condition theorem and restriction
to a direct summand.
Complex conjugation acts on the first block
by \(-1\) in the residual representation:
\(\bar\Psi(c_{\mathbb Q})=1\),
\(\bar\alpha(c_{\mathbb Q})=-1\).
Since \(\ell\) is odd and the matrix of complex
conjugation squares to one, its \(+1\)
eigenspace is zero (reduce its idempotent
projector). Hence its characteristic-zero
action is exactly \(-I\), and twisting by
\(\alpha^{-1}\) makes it \(I\).
This is evenness, not the oddness required
by the classical residually reducible ordinary
modularity theorem.

In characteristic \(\ell\), use the proper
flag specialization of
Theorem~\ref{thm:rh-sp4-curve-criterion}.
The two residual finite inertia characters
still belong to the first block by their
labeled actions at \(\mathrm{ph}\).
Its determinant is therefore everywhere of
finite-order inertia, and has finite global
order by class field theory.
An absolutely irreducible first block with
finite image would make both blocks have
finite image: the fixed similitude pairs the
blocks perfectly, so the second is the
dual of the first twisted by the residual
finite-order similitude.
The entire semisimplification would then have
finite image, giving trace equality and
diagonal type by the proof of
Theorem~\ref{thm:rh-sp4-curve-criterion},
a contradiction.
\end{proof}

\paragraph{Consequences for the source promodularity task.}
The characteristic-prime argument requires exclusion of
irreducible rank-four curves and irreducible rank-two
block endpoints, with the exact original local component.
Equivalently it is enough to prove finite-image
semisimplifications individually on all those curves.
A uniform finite quotient is unnecessary for this
part of the argument. The characteristic-zero endpoint
argument needs only (CF4) at one good prime on the
irreducible-block points. The irreducible rank-four
characteristic-zero attachment premise is the fixed-field
premise (A).

Neither exclusion in characteristic \(\ell\) has been
established here. The number-field group in this
theorem cannot be replaced by the fundamental group
of a curve over a finite field.
Nor can the even two-dimensional endpoint be declared
modular by ordinary lifting results which assume an
odd representation. The endpoint exclusion and the
irreducible attachment still require genuine arithmetic
input. The pinned source approaches these through
its rank-four promodularity tower, with a separate
mixed seed in Case II. In particular neither
characteristic-zero smoothness, the present support
classification, nor the finiteness of the base
rational cotangent space supplies those premises.
